\section{Chisel Compared with System Verilog}


%%%%%%%%%%%%%%%%%%%%%%
\subsection{Introduction}
%%%%%%%%%%%%%%%%%%%%%%


SystemVerilog and Chisel are both powerful languages used for hardware design and verification. SystemVerilog is an extension of Verilog, widely used in the industry for designing and verifying digital systems \cite{ieee1800-2017}. Chisel, on the other hand, is a hardware description language embedded in Scala, offering advanced features like parameterization and functional programming paradigms \cite{bachrach2012chisel}. This chapter provides a comparative analysis of SystemVerilog and Chisel, highlighting directly mapped constructs, unique features of each language, and guidance for implementing designs when transitioning between the two languages.

\subsection{Overview of SystemVerilog and Chisel}

\subsection{SystemVerilog}

SystemVerilog enhances Verilog by incorporating advanced modeling features, object-oriented programming concepts, and verification constructs \cite{baker2010systemverilog}. It is standardized as IEEE 1800 and is widely adopted in the semiconductor industry for both design and verification purposes.

\subsection{Chisel}

Chisel is a domain-specific language for hardware design embedded in Scala, enabling hardware designers to use modern programming paradigms \cite{chisel-doc}. It provides powerful abstractions for parameterization and reuse, and it generates synthesizable Verilog code, integrating with traditional hardware development flows.

\subsection{Directly Mapped Constructs}

Despite differences in syntax and paradigms, many hardware constructs have direct equivalents in both SystemVerilog and Chisel. This section presents a comparison of these constructs.

\subsection{Data Types}

\begin{table}[h!]
\centering
\begin{tabular}{|l|l|l|}
\hline
\textit{Concept} & \textit{SystemVerilog} & \textit{Chisel} \\
\hline
Boolean Signal & \texttt{logic} & \texttt{Bool} \\
Unsigned Integer & \texttt{logic [N-1:0]} & \texttt{UInt(N.W)} \\
Signed Integer & \texttt{signed logic [N-1:0]} & \texttt{SInt(N.W)} \\
Array & \texttt{logic [N-1:0] array [M-1:0]} & \texttt{Vec(M, UInt(N.W))} \\
Structure & \texttt{struct \{...\}} & \texttt{class ... extends Bundle \{...\}} \\
\hline
\end{tabular}
\caption{Directly Mapped Data Types}
\end{table}

\subsection{Modules and Hierarchy}

\begin{itemize}
    \item \textit{Module Definition}:
        \begin{itemize}
            \item \textbf{SystemVerilog}: Modules are defined using the \texttt{module} keyword.
            \item \textbf{Chisel}: Modules are classes extending \texttt{Module}.
        \end{itemize}
    \item \textit{Ports}:
        \begin{itemize}
            \item \textbf{SystemVerilog}: Ports are declared in the module header.
            \item \textbf{Chisel}: Ports are defined within an \texttt{IO} bundle.
        \end{itemize}
    \item \textit{Instantiation}:
        \begin{itemize}
            \item \textbf{SystemVerilog}: Instantiate modules by name and connect ports.
            \item \textbf{Chisel}: Instantiate modules by creating new instances and connecting \texttt{io} fields.
        \end{itemize}
\end{itemize}

\subsection{Assignments}

\begin{itemize}
    \item \textit{Continuous Assignment}:
        \begin{itemize}
            \item \textbf{SystemVerilog}: \texttt{assign out = a \& b;}
            \item \textbf{Chisel}: \texttt{out := a \& b}
        \end{itemize}
    \item \textit{Procedural Assignment}:
        \begin{itemize}
            \item \textbf{SystemVerilog}: Inside \texttt{always} blocks using \texttt{=} or \texttt{<=}
            \item \textbf{Chisel}: Use \texttt{:=} for signal assignment.
        \end{itemize}
\end{itemize}

\subsection{Control Structures}

\begin{itemize}
    \item \textit{Conditional Statements}:
        \begin{itemize}
            \item \textbf{SystemVerilog}: \texttt{if...else}
            \item \textbf{Chisel}: \texttt{when}, \texttt{.elsewhen}, \texttt{.otherwise}
        \end{itemize}
    \item \textit{Case Statements}:
        \begin{itemize}
            \item \textbf{SystemVerilog}: \texttt{case...endcase}
            \item \textbf{Chisel}: \texttt{switch}, \texttt{is}
        \end{itemize}
\end{itemize}

\subsection{Loops}

\begin{itemize}
    \item \textit{For Loops}:
        \begin{itemize}
            \item \textbf{SystemVerilog}: \texttt{for (int i = 0; i < N; i++) begin ... end}
            \item \textbf{Chisel}: \texttt{for (i <- 0 until N) \{ ... \}}
        \end{itemize}
    \item \textit{Generate Loops}:
        \begin{itemize}
            \item \textbf{SystemVerilog}: \texttt{generate...endgenerate}
            \item \textbf{Chisel}: Use Scala loops during elaboration to generate hardware.
        \end{itemize}
\end{itemize}

\subsection{Memory and Arrays}

\begin{itemize}
    \item \textit{Memories}:
        \begin{itemize}
            \item \textbf{SystemVerilog}: \texttt{logic [N-1:0] mem [0:M-1];}
            \item \textbf{Chisel}: \texttt{Mem(M, UInt(N.W))}
        \end{itemize}
\end{itemize}

\subsection{Features Unique to Chisel}

Chisel offers several features not available in SystemVerilog.

\subsection{Parameterization and Generics}

Chisel allows advanced parameterization using Scala's type system and functional programming features \cite{chisel-doc}.

\begin{itemize}
    \item \textit{Parameterized Modules}: Modules can be parameterized with types and values.
        \newline Example: \texttt{class MyModule(width: Int) extends Module \{...\}}
    \item \textit{Higher-Order Functions}: Functions that take other functions as parameters.
    \item \textit{Functional Programming Constructs}: Map, reduce, and other functional methods.
\end{itemize}

\subsection{Object-Oriented Programming}

Chisel leverages Scala's object-oriented features.

\begin{itemize}
    \item \textit{Inheritance}: Modules and Bundles can inherit from base classes.
    \item \textit{Traits}: Reusable collections of fields and methods.
\end{itemize}

\subsection{Type Safety and Strong Typing}

Chisel enforces strict type checking, reducing errors due to mismatched widths or types.

\subsection{Hardware Generation}

Chisel allows hardware generation using Scala code.

\begin{itemize}
    \item \textit{Dynamic Hardware Generation}: Create variable hardware structures based on parameters or conditions.
\end{itemize}

\subsection{Testing with ChiselTest}

Chisel provides a testing framework integrated with ScalaTest, enabling sophisticated testing methodologies.

\subsection{Features Unique to SystemVerilog}

SystemVerilog includes features not present in Chisel, particularly in verification.

\subsection{Verification Constructs}

SystemVerilog provides extensive verification features \cite{foster2003assertions}.

\begin{itemize}
    \item \textit{Assertions}: Immediate and concurrent assertions for design verification.
    \item \textit{Coverage}: Functional and code coverage constructs.
    \item \textit{Randomization}: Random stimulus generation with constraints.
\end{itemize}

\subsection{Interfaces and Modports}

SystemVerilog interfaces simplify module connections and support complex bus protocols \cite{guyer2005introduction}.

\subsection{Object-Oriented Verification}

SystemVerilog supports classes, inheritance, and polymorphism in testbenches.

\subsection{Time Control and Simulation}

SystemVerilog allows precise control over simulation time and events.

\subsection{Explicit Event Control}

SystemVerilog provides constructs like \texttt{\#delay}, \texttt{@(posedge clk)}, which are not directly available in Chisel.

\subsection{Implementing Designs: A Cross-Language Guide}

This section provides guidance for designers transitioning between SystemVerilog and Chisel.

\subsection{From SystemVerilog to Chisel}

\begin{itemize}
    \item \textit{Modules}: Replace \texttt{module} with classes extending \texttt{Module}.
    \item \textit{Ports}: Define an \texttt{IO} bundle instead of module port lists.
    \item \textit{Signals}: Use \texttt{val} and \texttt{Wire}, \texttt{Reg} instead of \texttt{wire}, \texttt{reg}.
    \item \textit{Assignments}: Use \texttt{:=} for signal assignments.
    \item \textit{Control Flow}: Replace \texttt{if...else} with \texttt{when...elsewhen...otherwise}.
    \item \textit{Generate Blocks}: Use Scala loops and conditionals for hardware generation.
\end{itemize}

\subsection{From Chisel to SystemVerilog}

\begin{itemize}
    \item \textit{Classes and Modules}: Map Chisel modules to \texttt{module} definitions.
    \item \textit{Bundles}: Translate \texttt{Bundle} classes to \texttt{struct} or module port groups.
    \item \textit{Functional Constructs}: Implement Scala functional constructs using procedural code or generate statements.
    \item \textit{Testing}: Replace ChiselTest with SystemVerilog testbenches and verification constructs.
\end{itemize}

\subsection{Common Pitfalls and Tips}

\begin{itemize}
    \item \textit{Type Differences}: Be mindful of type widths and signedness.
    \item \textit{Synchronous vs. Combinational Logic}: Ensure that registers and combinational logic are correctly implemented.
    \item \textit{Initialization}: SystemVerilog allows initial blocks, while Chisel uses \texttt{RegInit}.
    \item \textit{Simulation Time}: Chisel abstracts away simulation time; SystemVerilog allows explicit time control.
\end{itemize}




